\documentclass[12pt, a4paper]{article}

\usepackage[utf8]{inputenc}
\usepackage[T1]{fontenc}
\usepackage[italian, english]{babel}  % the last one is the main language
\usepackage{lmodern}
\usepackage[top=3cm, bottom=3cm]{geometry}
\usepackage{microtype}

\usepackage{amsmath, amssymb}
\usepackage{siunitx}
\usepackage{graphicx}
\usepackage{xcolor}
\usepackage{csquotes}

\usepackage[
  backend      = biber,
  style        = authoryear,
  maxcitenames = 2,
  maxbibnames  = 15,
  uniquelist   = false,
  uniquename   = false
]{biblatex}
\addbibresource{bibliography.bib}
\setlength{\bibitemsep}{0.5\baselineskip}

\usepackage[
  colorlinks = true,
  linkcolor  = black,
  citecolor  = blue!60!black,
  urlcolor   = blue!70!black
]{hyperref}

\begin{document}

\begin{center}
  {\large\bfseries Snow Water Equivalent in the Italian Alps: reconstruction from high-resolution 
  gridded temperature and precipitation datasets}
\end{center}

The seasonal snowpack acts as a natural water reservoir, storing solid winter precipitation and gradually releasing it as meltwater during the spring 
and summer months. In Northern Italy, snowmelt plays a crucial role in sustaining river discharge, accounting for approximately 50\% of the annual 
streamflow of the main rivers draining the Italian Alps \parencite{snow_drought}. However, snow cover is highly vulnerable to the ongoing global 
warming, as rising temperatures result in less snowfall and simultaneously enhance melt related processes: a widespread decline in snow depth (HS) 
has already been observed at in-situ station across the entire Alpine Arc \parencite{Matiu_HS}. Although station networks provide useful insights 
on snowpack dynamics, they are typically too sparse to accurately represent the marked spatial variability of snow cover. To this end, numerical 
modelling and remote sensing techniques have been developed to provide spatialized estimates. However, at the time of writing no available snow 
product allows for a reliable reconstruction of multi-decadal changes in snow water storage over Northern Italy. Indeed, dei tre dataset disponibili 
per le alpi italiane, ossia


By contrast, snow water equivalent (SWE), i.e. the 
amount of water stored within the snowpack, is measured far more sparsely, as its evaluation 



 Infatti, la misura dello snow water equivalent (SWE), ossia ..., 
viene fatta con cadenza nettamente inferiore rispetto a quella dello spessore nevoso, poichè richiede di scavare una trincea nivale. Questo fa sì 
che le uniche stime disponibili siano il risultato di modelli fisicamente basati, che simulano su delle griglie l'evoluzione del manto nevoso sotto 
meteorological forcing. Purtroppo per le alti italiane limitazioni principalmente legate all'ampiezza della finsestra temporale considerata rendono 
di poca utilità le due snow reanalisi incentrate sull'idrologia, IT-SNOW e Po River District, presenti per il nord italia.  

è time 
consuming perchè consiste nel Dato che misurare lo SWE, ossia ..., sono disponibile

Snow Water Equivalent (SWE)

the vast majority of available studies rely solely on manual station networks, which 
are too sparse to accurately represent the marked spatial variability of snow cover. At the time of writing, only three snow products 
providing spatialized estimates are available for the Italian Alps, namely MODIS, IT-SNOW and the Po River District dataset. However, any of these 
gridded datasets is suitable for reliably assessing the changes in 


con le perdite maggiori nei mesi primaverili. Furthermore, snow droghts, i.e. periods of 
persistently below average snowpack, have become both more frequent and longer-lasting since the 1990s. 

Alcuni effetti di 

In the last few decades, Europe has emerged as a hotspot for snow droughts, as periods of persistently below average snowpack have become both more 
frequent and longer-lasting since the 1990s. This trend appears to have intensified recently, with particularly severe snow deficits affecting 
European mountain ranges during the 2020s. In the Italian Alps, for instance, snow water storage reached record lows during both the 2022 and 2023 
winters, substantially altering the hydrological regime of Northern Italy as a consequence of reduced snowmelt runoff \parencite{snow_drought}.

Therefore, this thesis aims to address this gap by developing a modelling chain capable of simulating the seasonal snowpack between 1951 and 2023, 
providing new insights in snow cover changes across the Italian Hydrological portion of the Greater Alpine Range (IH-GAR).


The model presented in this thesis adopts a substantially minimalist approach, since a limited number of meteorological input is available over a 
period as long as 1951--2023. Indeed, only precipitation and temperature are required by our degree-day modelling chain.

The degree-day model was calibrated against in-situ snow depth records 

This study reveal a widespread
\printbibliography

\end{document}